\documentclass{article}

\usepackage{../handouts}

\title{CSCI 535 Handout 04: Homology}
\date{15 September 2026}
\author{Name: \rule{3in}{0.15mm}}

\begin{document}
\maketitle

\section*{Recall ...}

\begin{enumerate}
    \item If $(G,\heartsuit)$ is a group, what are the four properties that it satisfies?

        \textbf{Answer}
        \vspace{2in}

    \item If $(F,+,\cdot)$ is a field, what properties does it satisfy?

        \textbf{Answer}
        \vspace{2in}

    \pagebreak
    \begin{figure}[t!hb]
        \centering
        \begin{subfigure}[b]{0.3\textwidth}
            \includegraphics[width=\textwidth]{complex-1}
            \caption{Complex $J$}
            \label{fig:gull}
        \end{subfigure}
        ~ %
        \begin{subfigure}[b]{0.3\textwidth}
            \includegraphics[width=\textwidth]{complex-2}
            \caption{Complex $K$}
            \label{fig:tiger}
        \end{subfigure}
        ~ %
        \begin{subfigure}[b]{0.3\textwidth}
            \includegraphics[width=\textwidth]{complex-3}
            \caption{Complex $L$}
        \end{subfigure}
        \caption{Three complexes.}
    \end{figure}

    \item (Each working group is assigned one of the complexes above). For your complex,
        what are the chain groups, when thinking of the groups as sets?

        \textbf{Answer}
        \vspace{2in}

    \item What is a generating set of these groups? (A set $S \subseteq G$ is a
        \emph{generator set} if every element of $G$ can be constructed from
        elements of $S$ and their inverses, using the group operation).

        \textbf{Answer}
        \vspace{2in}

    \item Using $\Z \mod 3$ coefficients, what are the boundary operators?

        \textbf{Answer}
        \vspace{2in}

\end{enumerate}

\pagebreak
\section*{Quotient Groups}

If $(G,\diamondsuit)$ is a group and $N \subset of G$, we say that $N$ is a
\emph{subgroup} of $G$ if $N$ is also a group.  We say that $N$ is a
\emph{normal subgroup} if $N$ is invariant under conjugation by elements of $G$,
that is, for every $g \in G$, $g \diamondsuit N \diamondsuit g^{-1}$ is an
element of $N$.


\begin{enumerate}
    \item Every group has two normal subgroups.  What are they?

        \textbf{Answer}
        \vspace{2in}


    \item Give an example of a subgroup that is not a normal subgroup.

        \textbf{Answer}
        \vspace{2in}


    \pagebreak
    \item Consider $3\Z := \{3z ~|~ z \in \Z\}$.  Is $3\Z \triangleleft \Z$? Why
        or why not?

        \textbf{Answer}
        \vspace{3in}


    \item A \emph{coset} is Given a group $(G,\cdot)$, a subgroup $H$, and an element $g \in
        G$, the left coset is written as $gH:= \{g \cdot h ~|~ h \in H\}$. Given a
        a normal subgroup $N$ of $G$, the \emph{quotient group} $G/H := \{gN ~|~
        g \in G\}$. What is the quotient group $\Z / 3\Z$?

        \textbf{Answer}
        \vspace{2in}



\end{enumerate}

\pagebreak
\section*{Homology}

Given an $n \in \Z_+$, the homology group $H_n$ is defined as the quotient
group: $H_n := \ker(\partial_n)/\im(\partial_n)$.

\begin{enumerate}
    \item Using $\Z_3$ coefficients, what are the homology groups of your
        working group's simplicial complex?

        \textbf{Answer}
        \vspace{2in}


\end{enumerate}



\end{document}
